% ------------------------------------------------------------------------------
% 3 Methods
% ------------------------------------------------------------------------------
\section{Methods}\label{sec:methods}

\subsection{Frequentist efficacy analysis}\label{sec:methods-freq}

Disease-free survival and overall survival were summarized by
Kaplan--Meier curves with pointwise 95\% confidence intervals
(log--log transformation) and compared with log-rank tests, both overall
(2 degrees of freedom) and pairwise between arms adjusted for
multiplicity by closure of the testing family. Treatment effects were
estimated with Cox proportional-hazards models fitted by Efron's method
for ties. Three specifications are reported for the primary endpoint: an
unadjusted model; a covariate-adjusted complete-case model; and the
pre-specified covariate-adjusted model fitted on each of the $m=20$
imputed datasets and pooled by Rubin's rules~\cite{rubin1987}, with
degrees of freedom adjusted for the fraction of missing information.
The adjustment set comprised treatment arm, age, sex, obstruction,
perforation, adherence, number of positive nodes, differentiation, extent
of local spread, and surgery-to-registration interval. An analogous
pooled model was fitted for overall survival.

The proportional-hazards assumption was examined with scaled Schoenfeld
residuals~\cite{grambsch1994}, both graphically (loess smooths against
time) and with the global chi-square test. Two pre-specified sensitivity
analyses address potential violations and endpoint-definition concerns:
first, flexible functional forms for age and nodes via restricted
penalized splines (4 degrees of freedom each), compared with the linear
model by likelihood-ratio test; second, competing-risk analyses on the
recurrence endpoint, contrasting the Aalen--Johansen cumulative incidence
function with the naive one-minus-Kaplan--Meier estimate, and fitting
Fine--Gray subdistribution-hazard models~\cite{finegray1999} adjusted for
the same covariate set.

Because hazard ratios are not directly interpretable as absolute
benefit, five-year Kaplan--Meier survival probabilities with 95\%
confidence intervals were computed by arm for both endpoints, along with
between-arm risk differences and numbers needed to treat (NNT) whose
95\% confidence intervals were obtained by nonparametric bootstrap with
2{,}000 resamples (resampling patients, recomputing the Kaplan--Meier
estimates each time).

\subsection{Subgroup and interaction analyses}

Seven clinically motivated subgroups were pre-specified for the primary
comparison of levamisole plus 5-FU against observation: age
($<60$ vs.\ $\geqslant 60$ years), sex, nodal burden ($<4$ vs.\ $\geqslant 4$
positive nodes), obstruction, local extent (submucosa/muscle vs.\
serosa/contiguous), differentiation (well/moderate vs.\ poor), and
surgery-to-registration interval. Within each subgroup level, unadjusted
Cox models estimated subgroup-specific hazard ratios; effect modification
was tested by Wald tests of treatment-by-subgroup interaction terms in
models adjusted for the nodal count. Results are displayed as a forest
plot with interaction $p$-values. Given seven interactions tested, only
$p$-values below a Bonferroni-adjusted threshold of $0.05/7 \approx
0.0071$ would be interpreted as evidence of modification; this
multiplicity-aware reading is emphasized in
Section~\ref{sec:res-design}.

\subsection{Trial redesign and re-enacted monitoring}\label{sec:methods-design}

To redesign the trial under modern standards, planning assumptions were
taken from the re-analysis itself. The control arm's five-year DFS of
42.4\% and median DFS of 2.96 years imply an exponential baseline hazard
$\lambda_0 = \log 2 / 2.96 = 0.234$ per year under the (approximate)
constant-hazard reading; the clinically meaningful target was set at a
hazard ratio of $0.65$, close to the effect actually observed. A fixed
one-sided superiority design at $\alpha = 0.025$ with 90\% power,
1:1 allocation, three years of uniform accrual and a minimum of five
years of additional follow-up, requires 226 DFS events and approximately
322 patients (Schoenfeld's formula~\cite{schoenfeld1983}, verified with
\texttt{gsDesign::nSurv}~\cite{gsdesign2026}).

The modern redesign is a three-look group-sequential trial with analyses
after 50\% and 75\% of the planned events and at completion. Efficacy
boundaries use O'Brien--Fleming alpha-spending~\cite{o1979,lan-demets1983};
futility boundaries use a non-binding Hwang--Shih--DeCani beta-spending
with shape parameter $-4$~\cite{hwang1990}, which yields aggressive
early stopping for futility under the null while allowing the sponsor
to disregard the futility boundary without inflating the type-I error. The resulting
maximum sample size is 235 events (approximately 336 patients), a
4.1\% inflation over the fixed design---the standard price of
interim looks.

Two complementary analyses probe the design's behaviour. First, a power
curve displays the required events and achieved power as a function of
the target hazard ratio. Second, and more unusually, the monitoring plan
was \emph{re-enacted on the real trial data}: because the public dataset
contains no calendar dates, an information-time snapshot of the interim
stage was reconstructed by administratively censoring every patient's
follow-up at the time (years since randomization) at which the 118th
DFS event of the actual trial was observed, refitting the Cox
model on the censored-at-interim snapshot, and comparing the resulting
$Z$-statistic with the design boundaries. Conditional power for the final
analysis was computed under three assumptions for the remaining drift
(null effect, design alternative, and the observed interim trend), with
the sign convention that positive $Z$ favours the experimental arm.

\subsection{Simulation studies}\label{sec:methods-sim}

Three Monte Carlo studies quantify the operating characteristics behind
the design and analysis choices. The data-generating mechanism was
calibrated to the trial itself: covariates were drawn by bootstrap
resampling of the complete-case patients, the prognostic index used
hazard-ratio coefficients from the adjusted Cox model fit of
Section~\ref{sec:res-tx}, event times were exponential with hazard
$\lambda_0 \exp(\mathrm{lp}_i + \log(\mathrm{HR}) \cdot
\mathrm{treat}_i)$, accrual times were uniform on three years, and
administrative censoring occurred at eight years.

\textbf{Study A (covariate-adjustment efficiency)} simulated 1{,}000
trials per scenario for a two-arm comparison with $n=322$ at hazard
ratios of $1.0$ (type-I error) and $0.72$ (power), crossing three levels
of prognostic strength obtained by scaling the calibrated covariate
coefficients by 0, 0.5, and 1 (the corresponding standard deviations of
the prognostic index are 0, 0.208, and 0.416 on the log-hazard scale).
Each replicate fitted both the unadjusted Cox model and the
fully-adjusted model, recording estimates, model-based standard errors,
empirical standard deviations, and rejection frequencies, along with the
log-rank statistic. Because the hazard ratio is
non-collapsible~\cite{martinussen2013}, the unadjusted model targets the
marginal hazard ratio, which is attenuated toward the null relative to
the conditional (adjusted) estimand; the simulation quantifies both the
power gain from adjustment and this attenuation explicitly.

\textbf{Study B (group-sequential operating characteristics)} simulated
3{,}000 trials per scenario under the design of
Section~\ref{sec:methods-design} ($n = 335$, event triggers at
118/177/235). Interim snapshots were constructed by cutting each
simulated trial at the calendar time of the planned event count;
monitoring decisions followed the design boundaries, with futility
either adhered to or ignored (the non-binding option), giving three
scenarios: alternative with adherence, null with adherence, and null
with futility ignored. The study reports stopping probabilities by
analysis and reason, expected event counts, and the realized type-I
error and power, for direct comparison with the design-theoretic values
from \texttt{gsDesign}.

\textbf{Study C (interval estimation under non-proportional hazards)}
simulated 500 trials of $n=335$ in which the treatment effect was
deliberately time-varying: hazard ratio $0.45$ during the first year and
$0.80$ thereafter, a piecewise-constant pattern implemented through the
memoryless property of the exponential distribution. Each replicate
fitted an unadjusted Cox model and recorded the Wald 95\% confidence
interval for the log hazard ratio, together with a 100-resample
percentile bootstrap interval. Coverage was evaluated against the
simulation estimand (the mean of the 500 point estimates, i.e.\ the
quantity the estimator actually targets under misspecification), which
is the honest comparison for studying interval reliability.

\subsection{Machine-learning prediction}\label{sec:methods-ml}

Prognostic modelling treats the composite DFS endpoint as a prediction
problem, \emph{deliberately excluding the randomized treatment indicator
from all predictors} so that models quantify baseline risk only---a
clean separation between the prediction task of this section and the
causal tasks of Sections~\ref{sec:methods-bayes}
and~\ref{sec:methods-cate}. Five models were compared on the 888
complete cases: (1) a clinical Cox model with the same main effects as
the efficacy analysis; (2) an elastic-net penalized Cox model
($\ell_1$ ratio $0.5$, 50 alphas along the regularization path);
(3) a random survival forest (400 trees, minimum leaf size 8, $\sqrt{p}$
features per split); (4) gradient boosting with the Cox partial-likelihood
loss (250 trees, learning rate 0.05, minimum leaf size 8); and (5)
XGBoost with the Cox objective (300 rounds, depth 3, 80\% feature and
row subsampling).

Honest validation used $5\times5$-fold stratified cross-validation
(five repeats of five folds, stratified on the event indicator; 25
training--test partitions per model). Discrimination was measured by
Harrell's concordance index and by the inverse-probability-of-censoring
weighted (IPCW) AUC at five years; overall prediction error by the
integrated Brier score over years 1--5, using the Kaplan--Meier censoring
distribution for IPCW weighting~\cite{graf1999}. In-sample (apparent)
concordance is reported alongside the cross-validated values to expose
the optimism gap. The best model by cross-validated concordance was then
examined in depth: calibration by comparing mean predicted five-year
risks with Kaplan--Meier-observed risks within deciles of predicted
risk, a five-year time-ROC curve, and permutation importance (20
shuffles per covariate on a held-out 25\% test split).

\subsection{Bayesian re-analysis}\label{sec:methods-bayes}

The Bayesian analysis models the two-arm comparison (levamisole plus
5-FU vs.\ observation; $n = 607$ complete cases; 317 DFS events) with a
Weibull proportional-hazards regression fitted by Hamiltonian Monte
Carlo (No-U-Turn Sampler) in PyMC~\cite{salvatier2016}, with right censoring
handled by the \texttt{Censored} likelihood wrapper. Writing the
cumulative hazard as
\begin{equation}
H_i(t \mid x_i) = \left(\frac{t}{\beta_0 \exp(-\eta_i/\alpha)}\right)^{\alpha},
\qquad
\eta_i = \beta_{\mathrm{tr}}\,\mathrm{treat}_i + x_i^{\top}\gamma,
\label{eq:weibull}
\end{equation}
the shape parameter $\alpha$ permits increasing ($\alpha>1$) or
decreasing ($\alpha<1$) baseline hazards, and $\exp(\beta_{\mathrm{tr}})$
is the treatment hazard ratio. The covariate vector $x_i$ contains
age (per decade), $\log(1+\text{nodes})$, female sex, obstruction,
advanced extent, and long surgery-to-registration interval; priors for
covariate effects were weakly informative normals.

The centrepiece is a prior sensitivity analysis for the treatment
effect on the log-hazard-ratio scale, comparing three specifications
that span the spectrum of prior conviction: a \emph{skeptical} prior
$\mathcal{N}(0, 0.421^2)$, chosen so that only 5\% of the prior mass
lies below a hazard ratio of $0.5$; a \emph{weakly-informative}
$\mathcal{N}(0, 1)$; and an \emph{enthusiastic} prior
$\mathcal{N}(\log 0.65, 0.421^2)$ centred at the design alternative.
Posterior summaries (means, 95\% credible intervals, and the
probabilities $\Pr(\mathrm{HR}<1)$ and $\Pr(\mathrm{HR}<0.8)$) are
contrasted with each other and with the frequentist estimate.

Bayesian sequential monitoring refits the model under the skeptical
prior at four event-driven looks (25\%, 50\%, 75\%, and 100\% of the
observed DFS events), each time censoring the data at the time of the
corresponding event, exactly paralleling the frequentist re-enactment of
Section~\ref{sec:methods-design} so that the two monitoring philosophies
can be compared directly on the same evidence stream. Model criticism
comprises (i) a posterior predictive check replicating Kaplan--Meier
curves for the control arm from posterior draws, and (ii) approximate
leave-one-out cross-validation~\cite{vehtari2017} comparing the Weibull
baseline against an exponential special case ($\alpha \equiv 1$).

\subsection{Heterogeneous treatment effects}\label{sec:methods-cate}

The final module estimates average and conditional treatment effects on
the \emph{absolute} five-year DFS event risk for the two-arm comparison.
Jackknife pseudo-observations~\cite{andersen2003,parner2010} convert the
censored survival outcome into a continuous, regression-ready outcome:
for patient $i$,
\begin{equation}
\theta_i = n\,\hat\theta - (n-1)\,\hat\theta_{(-i)},
\end{equation}
where $\hat\theta$ is the Kaplan--Meier five-year event probability in
the full sample and $\hat\theta_{(-i)}$ the same quantity with patient
$i$ removed. On the modelling subset ($n = 594$ complete cases), the
population-average treatment effect was estimated three ways: the
unadjusted Kaplan--Meier risk difference with bootstrap confidence
intervals; cross-fitted G-computation~\cite{hernan2020} predicting both
arms' pseudo-observations with LASSO outcome models fitted in the
training folds; and a cross-fitted augmented inverse-probability
weighted (AIPW) estimator
\begin{equation}
\hat\psi_{\mathrm{AIPW}} = \frac{1}{n}\sum_{i=1}^{n}
\left\{\hat\mu_1(X_i) - \hat\mu_0(X_i)
+ \frac{T_i\,(\theta_i - \hat\mu_1(X_i))}{e}
- \frac{(1-T_i)\,(\theta_i - \hat\mu_0(X_i))}{1-e}\right\},
\label{eq:aipw}
\end{equation}
with influence-function standard errors, where
$\hat\mu_a(\cdot)$ are the cross-fitted arm-specific outcome models and
$e$ is the (known by design, empirically estimated) propensity
$P(T=1)=0.487$. Conditional effects were explored with a cross-fitted
T-learner (LASSO within each arm), yielding patient-level predicted
benefits $\hat\mu_0(X_i) - \hat\mu_1(X_i)$ whose distribution,
dependence on baseline risk, and quartile-level survival curves are
examined---with the explicit caveat that such exploratory
heterogeneity estimation in a single trial of this size is
hypothesis-generating only.

\subsection{Software and reproducibility}\label{sec:software}

All analyses were scripted in R 4.5.3 (packages \texttt{survival},
\texttt{mice}, \texttt{ggplot2}, \texttt{gsDesign}, and the
\texttt{tidyverse} stack) and Python 3.12 (\texttt{scikit-survival},
\texttt{lifelines}, \texttt{xgboost}, \texttt{PyMC}, \texttt{ArviZ});
exact versions are recorded in \texttt{environment/} within the
repository. A single global seed (20260916) plus module-specific seeds
govern every stochastic step, from imputation to posterior sampling.
The repository executes end-to-end via \texttt{make all}, regenerating
every figure and table in this report in approximately 20--25 minutes on
a single core. The pipeline has no hidden manual steps: data curation,
quality-control assertions (record counts, covariate consistency,
endpoint logic), analyses, and outputs are all code-reviewed and
deterministic.
